\documentclass{article}
\usepackage{graphicx}
\usepackage{geometry}
\usepackage[hidelinks]{hyperref}
\usepackage{parskip}
% \usepackage{biblatex}

\geometry{
    a4paper,
    total={170mm,257mm},
    left=20mm,
    top=20mm,
}

% \bibliography{references}


\begin{document}

\section{Introduction}

In Europe, cyling is, and in a growing number of regions, becoming a major alternative to motor vehicle use.

Primarly, single-cyclist crashes are events where no other road user is involved cite{}.

In the Netherlands, the risk of death in single-cyclist crashes increased by 7\% in the period from 1996 to 2014 cite{}

A single-cyclist crash is an interplay of three factors: vehicle, human, and environment cite{}.

Current knowledge leads to descriptive analysis of common factors and hypothesised mechanisms, such as slippery road conditions, visibility status, infrastructure concerns, among others cite{}.

To gather more accurate crash information, it is first necessary to objectively identify the individual crash features and cluster them by dynamic-related similarities.

The lack of objectivity generates bias in the assessment of crash causes cite{}.

Since single-cyclist crash classifications are neither standardised nor objectively assessed, we expose a knowledge gap regarding how crash information is gathered.


\section{Methods}

The approach to develop this classification consisted in three phases.

\subsection{Definitions}

\subsection{Literature review}

The search query for the literature on single-cyclist crashes was ...

Regarding the literature about bicycle dynamics, it is based on textbooks that cover the main areas of single-track vehicle dynamics.


\subsection{Single-cyclist crash anatomy}

A single-cyclist crash is the result of the rider losing control of the bicycle, which is the inability of the user to manoeuvre the vehicle to a desired state.

When task demand surpass a rider's capability, a critical scenario becomes a loss of control and a possible crash, depending on the extent of the loss of control.

When the critical scenario becomes a complete loss of control, it leads to a crash as defined in Section 2.1.

The pre-crash phase begins with a perturbation and ends when the rider no longer has the control authority to recover.

Then, the crash phase is composed of the uncontrolled motions as a result of the crash mechanisms, together with protective responses from the rider.

At last, the post-crash is typically brief and not well-defined in single-cyclist crashes, since due to the inherent instability of the bicycle, it will tend to rest immediately after the crash itself.


\subsection{Single-cyclist crash classifications}

Currently, studies on single-cyclist crashes are based on the available data, which is mainly descriptive from hospital, police, and insurance reports.

cite{} aimed to develop a crash detection system, which uses measured accelerations on the bicycle to detect the events.

Then, we find a four-layer approach from cite{}, using direct causes and fall types, which allows to analyse the kinematic form of cyclist’s falls. 

Although Gildea’s work analyses crash motions focused on the human, there is still a lack of information about the dynamics of the bicycle prior and leading to the crash.


\subsubsection{Known mechanisms}

Here we present the mechanisms already discussed in the literature.

Table

\subsection{Bicycle dynamics}

Bicycles are human-powered two-wheeled vehicles that require lateral balance control to maintain an upright position, which dynamics have been a matter of scientific interest for more than 100 years cite{}.

From the linear equations of motion, the stability of the bicycle is studied, showing a self-stable behaviour and three modes of vibration: weave, capsize, and caster.

Regarding the foundations of bicycles’ self-stability, cite{} demonstrated how it arises from several physics features of the bicycle.

For the human, the task of riding a bicycle consists in first maintaining balance and then, following a desired trajectory. 

Being that steering is the main control input and the fastest loop, through proprioceptive feedback, on the model of cite{}.

In addition, bicycles are sensitive to aerodynamic forces, which are studied as a resultant force acting on the centre of pressure.

Concerning bicycles under longitudinal acceleration, their limitations rely on load transfer and tyre grip cite{}.

Furthermore, tyre friction also plays a major role on cornering dynamics.

Tyre friction is a non-linear behaviour on lateral and longitudinal loads, which has a saturation limit at a given force.

Consequently, when the saturation limit is reached, the tyre starts to exhibit a phenomenon called ‘skidding’.

In summary, bicycles can be self-stable but require human control to maintain balance over a wider range of speeds. 


\subsection{Performance test of the proposed classification}

\subsubsection{Video dataset}

To test the proposed classification, we collected a set of single-cyclist crash videos in a dataset. 

\subsubsection{Questionnaire}

Following TU Delft’s Human Research Ethics Committee guidelines, we surveyed 20 participants, who received an informed consent which details the procedure of the experience, regarding data privacy, and participation.

The objective of the questionnaire was to test the clarity and scope of the classification.

The questionnaire was created using Google Forms, due to the flexibility available to make the questions.

\section{Results}

As shown in Section 2.3, single-cyclist crashes are the result of the rider losing control of the system due to an increase in the task demand. 

Crash causes have been assessed for more than 10 year.

\subsection{Sequential crash classification}

When the cause of accidents is assessed as a series of discrete events in a specific time, it is called ‘Linear chain-of-failure event model’ cite{}.

\subsubsection{Pre-crash manoeuvre}

In our classification framework, the first step consists of identifying the manoeuvre executed by the rider in the moment previous to the crash, which we cluster in the following options.

\subsubsection{Nature of the perturbation}

Due to the complexity of the human-bicycle system, it is subjected to several perturbations, here collected in three sub-classes.

\subsubsection{Mechanisms of crash}

In the literature, there are several known mechanisms of crashes, normally clustered as skidding, collisions, falls, among others, as seen in Table 1.

\subsubsection{Root causes}

The root causes serve a dual purpose, explaining the transition from the nature of the perturbation to the crash mechanisms, and as an open field for details about the factors that trigger the crash mechanism.

\subsubsection{Characteristic motion}

The last stage of a crash, which also causes more visual impact, is the motion of the system as a result of the crash mechanism.

From Section 2.5, we know that crash motions occur directly related to two of the angular degrees of freedom
of the rear frame of the bicycle, i.e. pitch and roll. 

\subsubsection{Unknown event features}

Although the research on bicycle and motorcycle dynamics has been extensive in the last 50 years, there are still major gaps to fill in the crash analysis field.

\subsubsection{Questionnaire results}

Each respondent of the questionnaire received a batch of ten videos, then, we used a score system to interpret the results.

\subsubsection{Classification diagram}

From the information gathered and explained above, we construct the diagram shown in Figure 7, which gives cues to follow a descriptive bicycle-crash classification. 

Special cases occur for minor and substantial loss of control.

\subsection{Single-cyclist crash classification example cases}

To illustrate the classification workflow, four representative single-cyclist crashes are presented in Figures 8, 9, 10, 11.

A low-side crash, potentially caused by the loss of grip on the rear tyre is shown in Figure 8. 

A high-side crash on a planned turning is shown in Figure 9.

The third case shows a rider on straight running, it is not possible to determine whether the rider is experiencing longitudinal acceleration or not (Figure 10).

The last case shows a rider manoeuvring through streets with tram tracks that need to be crossed in perpendicular directions (Figure 11).


\section{Discussion}

Here we propose a single-cyclist crash classification framework.

Possible failure mechanisms that lead to crashes were derived from single-track vehicle dynamics theory cite{}.

Our crash classification consists of five classes: pre-crash manoeuvre, nature of the perturbation, crash mechanisms, recoverability, and crash motion.

cite{}  includes the fall type in his classification, in our framework referred as crash motion.

\subsection{Crash mechanisms}

One important part of the present work is providing a detailed description of the hypothesised crash mechanisms.

Based on the crash model from cite{}, describing a crash mechanism as ‘loss of control’ becomes redundant, since all single-cyclist crashes occur due to loss of control

\subsection{Crash motion}

Since the most prominent characteristic of a crash is the motion, the ways of crashing are named after them, and the hypothesised mechanisms come as an additional description for the event.

\subsection{Questionnaire}

The questionnaire, designed to evaluate the scope of this study, revealed a high level of agreement regarding detected pre-crash manoeuvres, nature of the perturbation crash mechanisms, and motions.

\subsection{Edge cases for classification}

To show examples of the working of our classification, we chose four videos that could be challenging to classify with existing classifications.

Concerning example 3 (Figure 10), there is no evident infrastructure failure. 

At last, example 4 (Figure 11) would be classified as category 2-c-i in Schepers’ classification, i.e. loss of control due to abrupt steering manoeuvres.

\subsection{Limitations and future work}

The limitations of this work are related to the misunderstanding of some terms, such as the differences between tyre-ground interaction and collision forces. 

Altogether, this work contributes to the state of the art by providing a description of commonly used concepts that are not clearly defined in current articles in the field. 

\section{Conclusion}

From our review of the literature on single-cyclist crashes, we concluded that existing classification methods and mechanisms, do not adequately address bicycle dynamics

\end{document}